\documentclass[tikz,border=4pt]{standalone}
\usepackage{resistencia}
\begin{document}
% Sección del elipsoide de Lamé por un plano principal (σ1 y σ2): ejes principales como ejes.
\begin{tikzpicture}[x=1cm,y=1cm,etq/.style={etiq,fill=white,inner sep=1.2pt}]
\def\sa{3.4}\def\sb{1.8}
% circunferencia de normales unitarias y elipse de vectores tensión
\draw[gris,thick] (0,0) circle (1);
\draw[acero,very thick] (0,0) ellipse ({\sa} and {\sb});
\draw[eje] (-4.1,0) -- (4.3,0) node[right,etq]{$\sigma_1$};
\draw[eje] (0,-2.4) -- (0,2.6) node[above,etq]{$\sigma_2$};
% una normal n y su vector tensión t = (σ1 n1, σ2 n2)
\pgfmathsetmacro{\ang}{40}
\pgfmathsetmacro{\nx}{cos(\ang)}\pgfmathsetmacro{\ny}{sin(\ang)}
\pgfmathsetmacro{\tx}{\sa*\nx}\pgfmathsetmacro{\ty}{\sb*\ny}
\draw[vnormal] (0,0) -- (\nx,\ny) node[above left,etq,text=verde]{$\vec n$};
\draw[vtension,very thick] (0,0) -- (\tx,\ty) node[above right,etq]{$\vec t$};
\draw[dashed,tinta!60] (\tx,0) -- (\tx,\ty) -- (0,\ty);
\node[etq,text=acero] at (-2.2,2.25) {lugar de las puntas de $\vec t$};
\node[etq,text=tinta!70] at (-0.95,-0.55) {normales $\vec n$};
\fill (0,0) circle (1.4pt);
\end{tikzpicture}
\end{document}
